\chapter{The Advanced Encryption Standard (AES)}

In Chapter~11, DES and the Feistel architecture established the template for modern block ciphers, while showing why a 56-bit key and 64-bit block had to be replaced. When you open a web page over HTTPS, connect to a VPN, unlock an encrypted disk, or send a message on Signal today, the bytes leaving your machine are almost always scrambled by one replacement: the Advanced Encryption Standard (AES). Unlike DES, AES was designed in the open as a substitution--permutation network (SPN) on a 128-bit block, its every constant justified in public, and its mathematics is finite-field algebra over $\F_{256}$ that you can do by hand.

This chapter takes AES apart layer by layer: the four operations of a round (\texttt{SubBytes}, \texttt{ShiftRows}, \texttt{MixColumns}, \texttt{AddRoundKey}), a single round traced on real bytes, the S-box derived from the multiplicative inverse in $\F_{256}$, the key schedule, and the Wide Trail argument against differential and linear attacks.

\chapterline

\section{Where AES came from}

DES, the cipher of the 1970s and 1980s, was dying. Its $56$-bit key was small
enough that by the late 1990s a purpose-built machine (the Electronic Frontier
Foundation's ``Deep Crack'') could brute-force it in days. Triple-DES patched the
key size but tripled an already slow algorithm. The world needed a replacement,
and the U.S. National Institute of Standards and Technology (NIST) did something
unusual for a cryptographic standard: it held an open, international competition.

\begin{historybox}
In January 1997 NIST announced a call for a new block cipher to replace DES. The requirements: a $128$-bit block, support for $128$-, $192$-, and $256$-bit keys, royalty-free worldwide, and fast in both software and hardware. Fifteen submissions came from teams around the globe. Over three years cryptographers publicly attacked every candidate, and NIST held conferences to compare security, speed on everything from 8-bit smart cards to servers, and resistance to side-channel leakage. Five finalists emerged: MARS (IBM), RC6 (RSA Labs), Serpent (Anderson, Biham, Knudsen), Twofish (Schneier et al.), and Rijndael (Joan Daemen and Vincent Rijmen, two Belgians). Serpent was widely judged the most conservative (with 32 rounds and a huge security margin, though slower in software), while Rijndael used the Wide Trail strategy in 10 to 14 rounds to achieve provable bounds against differential and linear cryptanalysis with high speed across platforms. In October 2000 NIST chose Rijndael, and in 2001 it became FIPS 197: the Advanced Encryption Standard. The name ``Rijndael'' is a portmanteau of the designers' surnames; ``AES'' is the standardized profile of it (fixed $128$-bit block, three key sizes).
\end{historybox}

It is worth dwelling on what the open process bought us. DES's S-boxes looked
arbitrary, and for years people suspected a backdoor; only in the 1990s was it
revealed that IBM and the NSA had \emph{strengthened} them against differential
cryptanalysis, a technique the public would not rediscover until 1990. AES took
no such leap of faith. Every constant --- the irreducible polynomial, the S-box
affine map, the MixColumns matrix, the round constants --- is published with a
design rationale tied to measurable security goals. You can recompute all of them
yourself. That transparency is a big part of why the cryptographic community
trusts AES.

\begin{focusbox}
\textbf{The setting.} AES is a \emph{block cipher}: a keyed, invertible
permutation on $128$-bit blocks. Fix a key $k$ of $128$, $192$, or $256$ bits and
AES becomes a bijection $\Enc_k : \{0,1\}^{128} \to \{0,1\}^{128}$ with inverse
$\Dec_k$. By itself it encrypts exactly one $128$-bit block. To encrypt a message
of any length --- and to do so securely --- you wrap it in a \emph{mode of
operation} (CTR, CBC, GCM), the subject of the next chapter. Here we study the
permutation itself: what $\Enc_k$ does to one block.
\end{focusbox}

\section{The state: bytes in a $4\times 4$ grid}

AES works on bytes, not bits. A $128$-bit block is $16$ bytes, and AES arranges
them in a $4\times 4$ matrix called the \textbf{state}. The byte order is
\emph{column-major}: the first four input bytes fill the first column top to
bottom, the next four fill the second column, and so on.

If the input block is the byte sequence $b_0, b_1, \dots, b_{15}$, the state is

\[
\begin{bmatrix}
b_0 & b_4 & b_8 & b_{12}\\
b_1 & b_5 & b_9 & b_{13}\\
b_2 & b_6 & b_{10} & b_{14}\\
b_3 & b_7 & b_{11} & b_{15}
\end{bmatrix}.
\]

Every operation in AES is described as an action on this grid. \texttt{SubBytes}
replaces each cell independently. \texttt{ShiftRows} permutes within rows.
\texttt{MixColumns} mixes within columns. \texttt{AddRoundKey} XORs in key
material cell by cell. The grid shape matters for the diffusion argument
(Section~\ref{ch12:security}): ShiftRows
moves bytes \emph{between} columns and MixColumns spreads \emph{within} a column,
so the two together touch the whole block in a couple of rounds.

\begin{figure}[htbp]
\centering
\begin{tikzpicture}[
  >=Stealth,
  node distance=4mm,
  cell/.style={draw=gray!50, minimum size=4.6mm, font=\tiny\ttfamily, inner sep=0pt, rounded corners=0.5pt},
  arr/.style={->, thick, sourceblue!80!black},
  hdr/.style={font=\scriptsize\scshape, text=gray!80!black}
]
  \newcommand{\drawgrid}[6]{
    \begin{scope}[shift={#1}]
      \node[hdr] at (0.72, 0.72) {#2};
      \foreach \r/\cfill in {0/#3, 1/#4, 2/#5, 3/#6} {
        \foreach \c in {0,1,2,3} {
          \node[cell, fill=\cfill] at (\c*0.48, -\r*0.48) {};
        }
      }
    \end{scope}
  }

  % Stage 1: Input -> SubBytes -> ShiftRows
  \node[font=\footnotesize\scshape\bfseries, text=sourceblue!90!black] at (-1.3, -0.72) {Stage 1:};
  \drawgrid{(0,0)}{Initial State}{sourceblue!15}{targetgreen!15}{quotientorange!15}{structurepurple!15}

  \draw[arr] (1.9, -0.72) -- node[above, font=\scriptsize, align=center] {\texttt{SubBytes}\\(invert + affine)} (3.9, -0.72);
  \drawgrid{(4.2,0)}{\texttt{SubBytes}}{sourceblue!28}{targetgreen!28}{quotientorange!28}{structurepurple!28}

  \draw[arr] (6.1, -0.72) -- node[above, font=\scriptsize, align=center] {\texttt{ShiftRows}\\(row rotations)} (8.1, -0.72);

  % Grid 2: ShiftRows
  \begin{scope}[shift={(8.4,0)}]
    \node[hdr] at (0.72, 0.72) {\texttt{ShiftRows}};
    \foreach \c in {0,1,2,3} \node[cell, fill=sourceblue!28] at (\c*0.48, 0) {};
    \foreach \c in {0,1,2,3} \node[cell, fill=targetgreen!28] at (\c*0.48, -0.48) {};
    \foreach \c in {0,1,2,3} \node[cell, fill=quotientorange!28] at (\c*0.48, -0.96) {};
    \foreach \c in {0,1,2,3} \node[cell, fill=structurepurple!28] at (\c*0.48, -1.44) {};
    % Small shift arrows on rows 1, 2, 3
    \draw[->, thin, targetgreen!70!black] (1.65, -0.48) -- (1.35, -0.48);
    \draw[->, thin, quotientorange!80!black] (1.65, -0.96) -- (1.20, -0.96);
    \draw[->, thin, structurepurple!80!black] (1.65, -1.44) -- (1.05, -1.44);
  \end{scope}

  % Wrap-around connector from Stage 1 output to Stage 2 input
  \draw[arr, dashed] (10.3, -0.72) -- (10.8, -0.72) -- (10.8, -2.1) -- (-0.45, -2.1) -- (-0.45, -4.12) -- (-0.28, -4.12);

  % Stage 2: MixColumns -> AddRoundKey -> Output State
  \begin{scope}[yshift=-3.4cm]
  \node[font=\footnotesize\scshape\bfseries, text=sourceblue!90!black] at (-1.3, -0.72) {Stage 2:};
  \drawgrid{(0,0)}{Input to \texttt{Mix}}{sourceblue!28}{targetgreen!28}{quotientorange!28}{structurepurple!28}

  \draw[arr] (1.9, -0.72) -- node[above, font=\scriptsize, align=center] {\texttt{MixColumns}\\(MDS over $\F_{2^8}$)} (3.9, -0.72);
  \drawgrid{(4.2,0)}{\texttt{MixColumns}}{aliceblue!25}{aliceblue!25}{aliceblue!25}{aliceblue!25}

  \draw[arr] (6.1, -0.72) -- node[above, font=\scriptsize, align=center] {\texttt{AddRoundKey}\\($\xor \text{ Key}_r$)} (8.1, -0.72);
  \drawgrid{(8.4,0)}{Round Output}{targetgreen!30}{targetgreen!30}{targetgreen!30}{targetgreen!30}
  \end{scope}
\end{tikzpicture}
\caption{The four-layer state transformation across one round of AES. \texttt{SubBytes} applies nonlinear inversion and an affine map byte-by-byte. \texttt{ShiftRows} shifts rows cyclically to move data across columns. \texttt{MixColumns} mixes bytes within each column by matrix multiplication over $\F_{2^8}$. Finally, \texttt{AddRoundKey} injects the round subkey via XOR.}
\label{ch12:fig-aes-state-pipeline}
\end{figure}

\begin{intuitionbox}
Picture the state as a $4\times4$ checkerboard of bytes. AES alternates two kinds
of stirring. One kind scrambles each square in place (confusion). The other kind
slides squares around and blends neighbors so that a single changed input square
quickly contaminates the whole board (diffusion). Repeat enough times and the
output looks like noise, with no usable statistical handle for an attacker. The
rest of the chapter makes each of these steps precise.
\end{intuitionbox}

\subsection{Number of rounds}

AES is a \emph{substitution-permutation network} (SPN): there is no Feistel
split into halves as in DES. Instead every byte is transformed in every round.
The number of rounds depends on the key size:

\begin{center}
\begin{tabular}{lcc}
\toprule
Key size & Rounds $N_r$ & Round-key words \\
\midrule
AES-128 & 10 & 44 \\
AES-192 & 12 & 52 \\
AES-256 & 14 & 60 \\
\bottomrule
\end{tabular}
\end{center}

The high-level structure for AES-128 is:

\begin{enumerate}
\item \texttt{AddRoundKey} with round key $0$ (a ``whitening'' step).
\item Rounds $1$ through $9$, each: \texttt{SubBytes} $\to$ \texttt{ShiftRows}
$\to$ \texttt{MixColumns} $\to$ \texttt{AddRoundKey}.
\item Round $10$ (final): \texttt{SubBytes} $\to$ \texttt{ShiftRows} $\to$
\texttt{AddRoundKey}, with \textbf{no} MixColumns.
\end{enumerate}

The final round drops MixColumns for a clean reason we will see: it makes
decryption symmetric in structure with encryption and adds no security (the last
MixColumns could be undone by the attacker by linear algebra without the key, so
it would buy nothing).

\section{The field $\F_{256}$: arithmetic on bytes}
\label{ch12:field}

To understand \texttt{SubBytes} and \texttt{MixColumns} you need to think of a
byte as an element of the finite field $\F_{256} = \F_{2^8} = \mathrm{GF}(2^8)$.
We met finite fields in Chapter~6; here is the working summary AES needs.

A byte $b_7 b_6 \cdots b_1 b_0$ is read as the polynomial
\[
b_7 x^7 + b_6 x^6 + \cdots + b_1 x + b_0
\]
with coefficients in $\F_2 = \{0,1\}$. \textbf{Addition} of two bytes is just
polynomial addition mod $2$, i.e.\ bitwise XOR. There is no carry; addition and
subtraction are the same operation: $\mathtt{0x57} \xor \mathtt{0x83} = \mathtt{0xD4}$.

\textbf{Multiplication} is polynomial multiplication followed by reduction modulo
a fixed irreducible polynomial of degree $8$. AES fixes
\[
m(x) = x^8 + x^4 + x^3 + x + 1, \qquad \text{in hex, } \mathtt{0x11B}.
\]
This $m(x)$ has no factors over $\F_2$ (it is irreducible), which is exactly what
makes $\F_2[x]/(m(x))$ a field: every nonzero byte has a multiplicative inverse.
That is the property \texttt{SubBytes} exploits.

\begin{walkbox}
\textbf{Multiply $\mathtt{0x57}$ by $\mathtt{0x83}$ in $\F_{256}$.}

As polynomials, $\mathtt{0x57} = x^6 + x^4 + x^2 + x + 1$ and $\mathtt{0x83} =
x^7 + x + 1$. Multiply them out (over $\F_2$, so combine like terms mod $2$):
\[
(x^6 + x^4 + x^2 + x + 1)(x^7 + x + 1).
\]
Expanding term by term and cancelling pairs gives
\[
x^{13} + x^{11} + x^9 + x^8 + x^6 + x^5 + x^4 + x^3 + 1.
\]
This has degree $13 \ge 8$, so reduce modulo $m(x) = x^8+x^4+x^3+x+1$ using
$x^8 \equiv x^4+x^3+x+1 \pmod{m}$. First replace $x^{13} = x^5\cdot x^8$ by
$x^5(x^4+x^3+x+1) = x^9+x^8+x^6+x^5$: those four terms immediately cancel the
matching $x^9+x^8+x^6+x^5$ in the product, leaving $x^{11} + x^4 + x^3 + 1$.
Next replace $x^{11} = x^3\cdot x^8$ by $x^3(x^4+x^3+x+1) = x^7+x^6+x^4+x^3$:
now $x^4+x^3$ cancels in turn, leaving $x^7 + x^6 + 1 = \mathtt{0xC1}$. So
$\mathtt{0x57} \cdot \mathtt{0x83} = \mathtt{0xC1}$ (the canonical test vector
from the AES specification).
\end{walkbox}

In practice nobody multiplies two arbitrary bytes the long way. The trick AES
relies on is multiplication by $x$, called \texttt{xtime}, which is a left shift
plus a conditional XOR with $\mathtt{0x1B}$ (the low byte of $m(x)$):

\[
\texttt{xtime}(b) =
\begin{cases}
b \ll 1 & \text{if the top bit of } b \text{ is } 0,\\
(b \ll 1) \xor \mathtt{0x1B} & \text{if the top bit of } b \text{ is } 1.
\end{cases}
\]

Multiplying by $2$ is \texttt{xtime}; multiplying by $3$ is $\texttt{xtime}(b)
\xor b$; larger constants are built from these. MixColumns only ever multiplies
by $1$, $2$, and $3$, so \texttt{xtime} is all the field multiplication you need
for encryption.

\begin{connectionbox}
The companion algebra book builds $\F_{2^n}$ from scratch: why $\F_2[x]/(m(x))$
is a field exactly when $m$ is irreducible, how to find irreducible polynomials,
and why all fields of size $2^8$ are isomorphic (so AES's particular choice of
$m(x)$ is a convention, not a deep necessity --- a different irreducible
polynomial would give an equally good but incompatible cipher). The
multiplicative group $\F_{256}^{*}$ is cyclic of order $255$; that cyclic
structure is what guarantees inverses exist for the S-box.
\end{connectionbox}

\section{SubBytes: the only nonlinear layer}
\label{ch12:subbytes}

\texttt{SubBytes} replaces every byte of the state by another byte, using a
single fixed lookup table $S$ of $256$ entries (the \textbf{S-box}). It is the
sole source of nonlinearity in AES. Everything else --- ShiftRows, MixColumns,
AddRoundKey --- is linear (or affine) over $\F_2$, meaning it commutes with XOR.
If AES had no nonlinear step, the whole cipher would be one big affine map, and
an attacker could solve for the key with linear algebra. The S-box is what stops
that.

The S-box is built in two stages.

\textbf{Stage 1: the multiplicative inverse in $\F_{256}$.} Map each byte $a$ to
its inverse $a^{-1}$ in the field, with the convention that $\mathtt{0x00}$ maps
to itself (it has no inverse, so we just fix the missing point). This inverse map
is highly nonlinear --- it is, in a precise sense, about as far from linear as a
byte permutation can be, which is exactly why it resists differential and linear
cryptanalysis (Section~\ref{ch12:security}).

\textbf{Stage 2: an affine transformation over $\F_2$.} The pure inverse map has
one cosmetic weakness: it has fixed points (e.g.\ $0\mapsto 0$, $1\mapsto 1$) and
a clean algebraic description that an attacker might exploit. So AES post-composes
it with an invertible affine map. Writing the byte after inversion as bits
$(y_0,\dots,y_7)$, each output bit $s_i$ is

\[
s_i = y_i \xor y_{(i+4)\bmod 8} \xor y_{(i+5)\bmod 8} \xor y_{(i+6)\bmod 8}
\xor y_{(i+7)\bmod 8} \xor c_i,
\]

where the constant byte is $c = \mathtt{0x63} = (c_0,\dots,c_7) =
(1,1,0,0,0,1,1,0)$. In matrix form, over $\F_2$,

\[
\begin{bmatrix} s_0 \\ s_1 \\ s_2 \\ s_3 \\ s_4 \\ s_5 \\ s_6 \\ s_7 \end{bmatrix}
=
\begin{bmatrix}
1&0&0&0&1&1&1&1\\
1&1&0&0&0&1&1&1\\
1&1&1&0&0&0&1&1\\
1&1&1&1&0&0&0&1\\
1&1&1&1&1&0&0&0\\
0&1&1&1&1&1&0&0\\
0&0&1&1&1&1&1&0\\
0&0&0&1&1&1&1&1
\end{bmatrix}
\begin{bmatrix} y_0 \\ y_1 \\ y_2 \\ y_3 \\ y_4 \\ y_5 \\ y_6 \\ y_7 \end{bmatrix}
\xor
\begin{bmatrix} 1 \\ 1 \\ 0 \\ 0 \\ 0 \\ 1 \\ 1 \\ 0 \end{bmatrix}.
\]

The matrix is a \emph{circulant}: each row is the previous one rotated, which is
why the formula above is just ``XOR five rotated copies.''

\begin{whybox}
Why split SubBytes into inversion-then-affine? The inverse map $a \mapsto a^{-1}$
gives the cryptographic strength: it has the best known resistance to
differential attacks (its differential uniformity is $4$, near the theoretical
minimum of $2$) and to linear attacks. But it is too ``algebraic'' on its own ---
it is described by a simple equation over $\F_{256}$, and it fixes $0$ and $1$.
The affine layer is cheap insurance: it breaks the algebraic simplicity over the
field (the affine map is $\F_2$-linear but \emph{not} $\F_{256}$-linear, so the
composition has no clean field description), removes fixed points, and ensures
$S(a) \neq a$ and $S(a) \neq \bar a$ for all $a$. Because the affine map is
invertible and linear over $\F_2$, it does not weaken the differential or linear
properties that the inverse map provides --- it simply destroys polynomial tractability.
\end{whybox}

\begin{walkbox}
\textbf{Compute $S(\mathtt{0x53})$ from scratch.}

First invert $\mathtt{0x53}$ in $\F_{256}$. You can check (by the extended
Euclidean algorithm on polynomials, or a log table) that $\mathtt{0x53}^{-1} =
\mathtt{0xCA}$. Now apply the affine map to $y = \mathtt{0xCA} =
(y_0,\dots,y_7) = (0,1,0,1,0,0,1,1)$ (bit $0$ is the least significant). Computing
each $s_i = y_i \xor y_{i+4} \xor y_{i+5} \xor y_{i+6} \xor y_{i+7} \xor c_i$
(indices mod $8$) gives the byte $\mathtt{0xED}$. And indeed, looking up
$\mathtt{0x53}$ in the published S-box table returns $\mathtt{0xED}$. The standard
AES S-box lookup table simply records the output of this two-stage calculation
evaluated across all $256$ byte values and stored for fast execution.
\end{walkbox}

\begin{implementbox}
A toy S-box generator in Python. It builds the field inverse with a log/antilog
table over generator $3$, then applies the affine map. \textbf{For learning only}
--- real code uses the precomputed table (or hardware) and must avoid the
data-dependent memory lookups shown here, which leak timing (see
Chapter~37).

\begin{lstlisting}[language=Python]
def xtime(b):                 # multiply by x (i.e. by 2) in GF(2^8)
    b <<= 1
    if b & 0x100:             # overflow past degree 7?
        b ^= 0x11B            # reduce mod x^8+x^4+x^3+x+1
    return b & 0xFF

def gmul(a, b):               # full field multiply
    p = 0
    for _ in range(8):
        if b & 1:
            p ^= a
        a = xtime(a)
        b >>= 1
    return p

# build log / antilog tables using generator g = 3
antilog = [0] * 256
x = 1
for i in range(255):
    antilog[i] = x
    x = gmul(x, 3)            # 3 is a generator of GF(2^8)*
log = [0] * 256
for i in range(255):
    log[antilog[i]] = i

def inv(a):                   # multiplicative inverse, 0 -> 0
    if a == 0:
        return 0
    return antilog[(255 - log[a]) % 255]

def rotl8(v, n):
    return ((v << n) | (v >> (8 - n))) & 0xFF

def sbox(a):
    y = inv(a)
    # affine: s = y XOR (y<<<1) XOR (y<<<2) XOR (y<<<3) XOR (y<<<4) XOR 0x63
    s = y ^ rotl8(y, 1) ^ rotl8(y, 2) ^ rotl8(y, 3) ^ rotl8(y, 4) ^ 0x63
    return s & 0xFF

assert sbox(0x53) == 0xED
assert sbox(0x00) == 0x63     # 0 -> 0 under inverse, then +0x63
print([hex(sbox(i)) for i in range(8)])
\end{lstlisting}
\end{implementbox}

\section{ShiftRows: cheap cross-column movement}
\label{ch12:shiftrows}

\texttt{SubBytes} treats each byte independently --- it never mixes bytes
together. If that were the whole cipher, byte $0$ of the ciphertext would depend
only on byte $0$ of the plaintext, and AES would be $16$ tiny independent ciphers.
We need to spread information across positions. \texttt{ShiftRows} is the first,
cheapest spreading step: it cyclically rotates the rows of the state.

\[
\begin{bmatrix}
a_{0} & a_{4} & a_{8} & a_{12}\\
a_{1} & a_{5} & a_{9} & a_{13}\\
a_{2} & a_{6} & a_{10} & a_{14}\\
a_{3} & a_{7} & a_{11} & a_{15}
\end{bmatrix}
\;\xrightarrow{\ \texttt{ShiftRows}\ }\;
\begin{bmatrix}
a_{0} & a_{4} & a_{8} & a_{12}\\
a_{5} & a_{9} & a_{13} & a_{1}\\
a_{10} & a_{14} & a_{2} & a_{6}\\
a_{15} & a_{3} & a_{7} & a_{11}
\end{bmatrix}
\]

Row $0$ is left unchanged. Row $1$ is rotated left by one byte, row $2$ by two,
row $3$ by three. It is a pure permutation of byte positions --- no arithmetic,
no key, just relabeling. On its own it does nothing cryptographically; its job is
to set up MixColumns. Because MixColumns mixes \emph{within} a column, you need
something to move bytes \emph{between} columns first, or the four columns would
evolve as four independent sub-ciphers forever. ShiftRows guarantees that after
one round the four bytes that were in one column are spread across all four
columns, so the next MixColumns blends them with bytes from elsewhere.

\begin{intuitionbox}
ShiftRows and MixColumns work together. MixColumns mixes up-and-down inside a column;
ShiftRows mixes left-and-right across columns. Together, each output byte after
two rounds depends on \emph{every} input byte. Neither alone achieves this:
MixColumns keeps each column isolated, ShiftRows keeps each byte's value
untouched. Together they give AES full diffusion in two rounds.
\end{intuitionbox}

\section{MixColumns: diffusion inside a column}
\label{ch12:mixcolumns}

\texttt{MixColumns} is the heavy diffusion step. It treats each column of the
state as a vector of four bytes and multiplies it by a fixed $4\times4$ matrix
over $\F_{256}$:

\[
\begin{bmatrix} a_0' \\ a_1' \\ a_2' \\ a_3' \end{bmatrix}
=
\begin{bmatrix}
\mathtt{02} & \mathtt{03} & \mathtt{01} & \mathtt{01}\\
\mathtt{01} & \mathtt{02} & \mathtt{03} & \mathtt{01}\\
\mathtt{01} & \mathtt{01} & \mathtt{02} & \mathtt{03}\\
\mathtt{03} & \mathtt{01} & \mathtt{01} & \mathtt{02}
\end{bmatrix}
\begin{bmatrix} a_0 \\ a_1 \\ a_2 \\ a_3 \end{bmatrix},
\]

where all matrix operations are performed over the finite field $\F_{2^8}$. Recall how matrix-vector multiplication operates: the $i$-th element of the resulting vector is the inner product of the $i$-th row of the matrix with the input column vector:
\[
a_i' = \bigoplus_{j=0}^{3} \big(M_{i,j} \cdot a_j\big).
\]
In $\F_{2^8}$, the field addition $\bigoplus$ is bitwise XOR ($\xor$), and field multiplication $\cdot$ is polynomial multiplication modulo $m(x) = x^8 + x^4 + x^3 + x + 1$. Because every entry in the matrix is either $\mathtt{01}$, $\mathtt{02}$, or $\mathtt{03}$, computing each product requires no general polynomial long division:
\begin{itemize}
  \item $\mathtt{01} \cdot a = a$ (the multiplicative identity).
  \item $\mathtt{02} \cdot a = \texttt{xtime}(a)$: shift byte $a$ left by 1 bit, and if its high bit before shifting was 1, XOR with the reduction byte $\mathtt{1B}$.
  \item $\mathtt{03} \cdot a = (\mathtt{02} \xor \mathtt{01}) \cdot a = \texttt{xtime}(a) \xor a$.
\end{itemize}
Concretely, the first output byte is:
\[
a_0' = (\mathtt{02}\cdot a_0) \xor (\mathtt{03}\cdot a_1) \xor a_2 \xor a_3,
\]
and the remaining three rows are cyclic rotations of this exact coefficient pattern (Figure~\ref{ch12:mixcolumns-fig}). The matrix is circulant, derived from multiplying by the fixed polynomial $c(x) = \mathtt{03}\,x^3 + \mathtt{01}\,x^2 + \mathtt{01}\,x + \mathtt{02}$ modulo $x^4+1$ over $\F_{2^8}$.

\begin{walkbox}
\textbf{Mix one column.} Take the column $(\mathtt{D4}, \mathtt{BF},
\mathtt{5D}, \mathtt{30})$, a value that appears in the FIPS-197 worked example.
Compute the first output byte:
\[
a_0' = (\mathtt{02}\cdot \mathtt{D4}) \xor (\mathtt{03}\cdot \mathtt{BF})
\xor \mathtt{5D} \xor \mathtt{30}.
\]
Now $\mathtt{02}\cdot\mathtt{D4} = \texttt{xtime}(\mathtt{D4})$: the top bit of
$\mathtt{D4} = 11010100$ is set, so shift left to $\mathtt{1A8}$ and XOR
$\mathtt{1B}$, giving $\mathtt{A8}\xor\mathtt{1B}=\mathtt{B3}$. And
$\mathtt{03}\cdot\mathtt{BF} = \texttt{xtime}(\mathtt{BF}) \xor \mathtt{BF}$;
$\texttt{xtime}(\mathtt{BF})$: top bit of $\mathtt{BF}=10111111$ is set, shift to
$\mathtt{17E}$, XOR $\mathtt{1B}$ gives $\mathtt{65}$, then
$\mathtt{65}\xor\mathtt{BF} = \mathtt{DA}$. So
\[
a_0' = \mathtt{B3} \xor \mathtt{DA} \xor \mathtt{5D} \xor \mathtt{30} = \mathtt{04}.
\]
Working out all four rows the same way gives the output column $(\mathtt{04},
\mathtt{66}, \mathtt{81}, \mathtt{E5})$. The matrix has spread each input byte
into all four outputs.
\end{walkbox}

\begin{whybox}
Why \emph{these} constants ($\mathtt{02},\mathtt{03},\mathtt{01},\mathtt{01}$)?
Two reasons.

\textbf{Maximal diffusion.} The MixColumns matrix is the generating matrix of a
\emph{maximum-distance-separable} (MDS) code. The practical consequence is the
``branch number'' is $5$, the largest possible for a $4\times4$ matrix: if you
change $d$ input bytes of a column ($1\le d\le 4$), then the number of changed
input bytes plus changed output bytes is at least $5$. So flipping even a single
input byte forces all four output bytes to change. This is the strongest possible
single-column avalanche.

\textbf{Cheapness.} The constants $\mathtt{01},\mathtt{02},\mathtt{03}$ are the
smallest values that achieve MDS, and they need only \texttt{xtime} and XOR --- no
general multiplier. A designer who wanted MDS but used bigger constants would pay
for it in every implementation forever. Daemen and Rijmen searched for the
cheapest MDS matrix; this is it.
\end{whybox}

\begin{figure}[H]
\centering
\begin{tikzpicture}[font=\scriptsize, scale=0.9, every node/.style={transform shape}]
  % input column
  \foreach \i/\v in {0/$a_0$,1/$a_1$,2/$a_2$,3/$a_3$} {
    \node[draw, minimum size=0.7cm, fill=sourceblue!12] (in\i) at (0,-\i) {\v};
  }
  \node[above=0.15cm of in0] {input column};
  % matrix label
  \node (mat) at (2.6,-1.5) [draw, rounded corners, fill=targetgreen!12,
        minimum width=2.0cm, minimum height=2.4cm, align=center]
        {MDS\\ matrix\\ $\times$};
  % output column
  \foreach \i/\v in {0/$a_0'$,1/$a_1'$,2/$a_2'$,3/$a_3'$} {
    \node[draw, minimum size=0.7cm, fill=quotientorange!18] (out\i) at (5.2,-\i) {\v};
  }
  \node[above=0.15cm of out0] {output column};
  % arrows: every input touches every output (branch number 5)
  \foreach \i in {0,1,2,3} {
    \foreach \j in {0,1,2,3} {
      \draw[-{Stealth[length=1.4mm]}, gray!55] (in\i.east) -- (out\j.west);
    }
  }
\end{tikzpicture}
\caption{\texttt{MixColumns} multiplies each $4$-byte column by a fixed MDS
matrix over $\F_{256}$. Because the branch number is $5$, every input byte
influences every output byte (all sixteen arrows are live), so changing one input
byte changes all four outputs.}
\label{ch12:mixcolumns-fig}
\end{figure}

\begin{connectionbox}
The branch-number argument is coding theory: the rows of the MixColumns matrix,
appended to the identity, generate a $[8,4,5]$ MDS code over $\F_{256}$, whose
minimum distance $5$ \emph{is} the branch number. The companion algebra book's
treatment of Reed--Solomon and MDS codes is exactly the machinery that proves an
attacker cannot find a low-weight differential through MixColumns, which links
error-correcting codes to provable resistance to differential cryptanalysis.
\end{connectionbox}

\section{AddRoundKey: where the key finally enters}
\label{ch12:addroundkey}

The three layers so far are fixed and public --- they do not depend on the key at
all. \texttt{AddRoundKey} is the only step that uses secret material, and it could
not be simpler: XOR the $128$-bit round key into the $128$-bit state, byte for
byte (equivalently, column for column).

\[
\text{state} \;\gets\; \text{state} \;\xor\; \text{RoundKey}_r.
\]

That is the whole operation. Its own inverse is itself (XOR twice with the same
key returns the original). It looks trivial, but it is essential: without it, $\Enc_k$ would be a fixed public permutation with no secret, and
anyone could decrypt. With it, the attacker faces a different permutation for
every key, and the nonlinear SubBytes layers wrapped around the key additions
make it computationally infeasible to peel the keys back off.

\begin{trapbox}
A naive reaction is ``AddRoundKey is just XOR, that is weak, like a one-time pad
with a reused key.'' Two things save it. First, the key material is fresh per
round (the key schedule generates a new round key each time). Second, and more
importantly, AddRoundKey is sandwiched between \emph{nonlinear} SubBytes layers.
A plain XOR with a repeated key is breakable because the cipher is linear; AES is
not. The key XOR is the simplest possible way to inject the secret, and the
nonlinearity around it does the protecting.
\end{trapbox}

\section{The key schedule}
\label{ch12:keyschedule}

AES-128 needs eleven $128$-bit round keys (one for the initial whitening plus one
per round). The \textbf{key schedule} (or key expansion) stretches the
$128$-bit cipher key into these $11\times16 = 176$ bytes. It is organized as a
sequence of $4$-byte \emph{words} $W_0, W_1, \dots, W_{43}$. The first four words
are just the cipher key; each later word is derived from earlier ones.

For AES-128:

\begin{itemize}
\item $W_0,W_1,W_2,W_3$ = the four words of the cipher key.
\item For $i = 4,\dots,43$:
\[
W_i =
\begin{cases}
W_{i-4} \xor \big(\texttt{SubWord}(\texttt{RotWord}(W_{i-1})) \xor \texttt{Rcon}_{i/4}\big)
& \text{if } 4 \mid i,\\[4pt]
W_{i-4} \xor W_{i-1} & \text{otherwise.}
\end{cases}
\]
\end{itemize}

Here \texttt{RotWord} rotates a word's four bytes left by one ($[a_0,a_1,a_2,a_3]
\to [a_1,a_2,a_3,a_0]$), \texttt{SubWord} applies the S-box to each of the four
bytes, and the \textbf{round constant} $\texttt{Rcon}_j = (\mathrm{rc}_j,0,0,0)$
has $\mathrm{rc}_1 = \mathtt{01}$ and $\mathrm{rc}_{j} = \texttt{xtime}(\mathrm{rc}_{j-1})$,
i.e.\ the powers of $x$ in $\F_{256}$:
\[
\mathtt{01},\ \mathtt{02},\ \mathtt{04},\ \mathtt{08},\ \mathtt{10},\
\mathtt{20},\ \mathtt{40},\ \mathtt{80},\ \mathtt{1B},\ \mathtt{36}.
\]
Round key $r$ is the concatenation of words $W_{4r}, W_{4r+1}, W_{4r+2},
W_{4r+3}$, as shown in Figure~\ref{ch12:fig-key-schedule}.

\begin{figure}[htbp]
\centering
\begin{tikzpicture}[
  >=Stealth,
  word/.style={draw=sourceblue!80!black, fill=sourceblue!10, rounded corners=2pt, minimum width=1.4cm, minimum height=0.65cm, font=\scriptsize\bfseries},
  op/.style={draw=structurepurple!80!black, fill=structurepurple!12, rounded corners=2pt, minimum width=1.65cm, minimum height=0.52cm, font=\scriptsize\ttfamily},
  xorgate/.style={draw=quotientorange!90!black, fill=quotientorange!15, circle, inner sep=0pt, minimum size=4.6mm, font=\scriptsize\bfseries},
  arr/.style={->, thick, gray!80!black}
]
  % Top row: Round key r-1 words W0..W3
  \node[font=\scriptsize\scshape, text=gray!70!black, anchor=east] at (-0.9, 0) {Round Key $0$:};
  \node[word] (w0) at (0, 0) {$W_0$};
  \node[word] (w1) at (2.2, 0) {$W_1$};
  \node[word] (w2) at (4.4, 0) {$W_2$};
  \node[word] (w3) at (6.6, 0) {$W_3$};

  % g-function pipeline on the right for W3
  \node[op] (rot) at (9.5, -0.45) {RotWord};
  \node[op] (sub) at (9.5, -1.25) {SubWord};
  \node[xorgate] (xrcon) at (9.5, -2.05) {$\xor$};
  \node[draw=quotientorange!80!black, fill=quotientorange!10, rounded corners=1.5pt, font=\scriptsize\ttfamily, inner sep=3pt] (rcon) at (11.4, -2.05) {Rcon$_1$};

  \draw[arr, structurepurple!80!black] (w3.east) -| (rot.north);
  \draw[arr, structurepurple!80!black] (rot) -- (sub);
  \draw[arr, structurepurple!80!black] (sub) -- (xrcon);
  \draw[arr, quotientorange!80!black] (rcon) -- (xrcon);

  % Dashed box around g
  \draw[dashed, rounded corners=3pt, structurepurple!60] (8.45, -0.05) rectangle (12.25, -2.45);
  \node[font=\scriptsize\bfseries, text=structurepurple!85!black, anchor=south] at (10.35, -0.05) {nonlinear core $g(W_3)$};

  % XOR gates for W4..W7
  \node[xorgate] (x4) at (0, -2.8) {$\xor$};
  \node[xorgate] (x5) at (2.2, -2.8) {$\xor$};
  \node[xorgate] (x6) at (4.4, -2.8) {$\xor$};
  \node[xorgate] (x7) at (6.6, -2.8) {$\xor$};

  % Route g(W3) output back to x4
  \draw[arr, structurepurple!80!black] (xrcon.south) -- (9.5, -2.8) -- (7.4, -2.8) -- (7.4, -1.6) -- (-0.55, -1.6) |- (x4.west);

  % Vertical feed from W0..W3 to x4..x7
  \draw[arr] (w0) -- (x4);
  \draw[arr] (w1) -- (x5);
  \draw[arr] (w2) -- (x6);
  \draw[arr] (w3) -- (x7);

  % Bottom row: Round key 1 words W4..W7
  \node[font=\scriptsize\scshape, text=gray!70!black, anchor=east] at (-0.9, -3.85) {Round Key $1$:};
  \node[word, fill=targetgreen!14, draw=targetgreen!80!black] (w4) at (0, -3.85) {$W_4$};
  \node[word, fill=targetgreen!14, draw=targetgreen!80!black] (w5) at (2.2, -3.85) {$W_5$};
  \node[word, fill=targetgreen!14, draw=targetgreen!80!black] (w6) at (4.4, -3.85) {$W_6$};
  \node[word, fill=targetgreen!14, draw=targetgreen!80!black] (w7) at (6.6, -3.85) {$W_7$};

  \draw[arr] (x4) -- (w4);
  \draw[arr] (x5) -- (w5);
  \draw[arr] (x6) -- (w6);
  \draw[arr] (x7) -- (w7);

  % Chaining W4 -> x5, W5 -> x6, W6 -> x7
  \draw[arr, targetgreen!75!black] (0, -3.3) -- (1.1, -3.3) |- (x5.west);
  \draw[arr, targetgreen!75!black] (2.2, -3.3) -- (3.3, -3.3) |- (x6.west);
  \draw[arr, targetgreen!75!black] (4.4, -3.3) -- (5.5, -3.3) |- (x7.west);
\end{tikzpicture}
\caption{One step of the AES-128 key schedule expanding $(W_0, W_1, W_2, W_3)$ into $(W_4, W_5, W_6, W_7)$. Every fourth word ($4 \mid i$) passes the previous word through the nonlinear core $g$ (\texttt{RotWord}, \texttt{SubWord}, and $\xor\,\texttt{Rcon}_{i/4}$) before XORing with $W_{i-4}$; the remaining three words simply XOR $W_{i-1}$ with $W_{i-4}$.}
\label{ch12:fig-key-schedule}
\end{figure}

\begin{whybox}
Every step in the key schedule serves a distinct security purpose.

\textbf{Nonlinearity (\texttt{SubWord}).} Without the S-box in the schedule, the
round keys would be a linear function of the cipher key, and an attacker who
learned one round key could solve for the others. SubWord makes the schedule
nonlinear so round keys cannot be linearly related.

\textbf{Symmetry breaking (\texttt{Rcon}).} The round constants differ every round
(they are distinct powers of $x$). Without them, the schedule would treat each
round identically, creating symmetries an attacker could exploit (related-key and
slide attacks thrive on round-to-round symmetry). Distinct $\texttt{Rcon}$ values
make every round key structurally different.

\textbf{Rotation (\texttt{RotWord}).} Rotating before substitution diffuses byte
positions across the word, so a change in one key byte spreads to others.
\end{whybox}

\begin{walkbox}
\textbf{Expand the all-zero key, first step.} Let the cipher key be sixteen zero
bytes, so $W_0=W_1=W_2=W_3 = \mathtt{00000000}$. Compute $W_4$. Since $4\mid 4$:
\texttt{RotWord}($W_3$)$ = \mathtt{00000000}$, \texttt{SubWord} of that applies
$S(\mathtt{00})=\mathtt{63}$ to each byte giving $\mathtt{63636363}$, then XOR
$\texttt{Rcon}_1 = \mathtt{01000000}$ gives $\mathtt{62636363}$. Finally $W_4 =
W_0 \xor \mathtt{62636363} = \mathtt{62636363}$. Then $W_5 = W_1 \xor W_4 =
\mathtt{62636363}$, $W_6 = W_2 \xor W_5 = \mathtt{62636363}$, $W_7 = W_3 \xor W_6
= \mathtt{62636363}$. So round key $1$ is $\mathtt{62636363}$ repeated four times.
The S-box and Rcon have already injected structure into what was an all-zero key.
\end{walkbox}

\begin{trapbox}
The key schedule is the part of AES with the smallest security margin, and it is
where most academic ``breaks'' live. The known \emph{related-key} attacks on
AES-256 (Biryukov--Khovratovich, 2009) exploit the simpler key schedule of the
$256$-bit variant: with $14$ rounds but a schedule that does relatively little
mixing per round key, an attacker who can query the cipher under keys with chosen
\emph{differences} can recover the key faster than brute force. This is mostly of
theoretical concern --- related-key access is a strong and unusual assumption,
and no attack threatens AES under a single fixed unknown key --- but it shows
that the key schedule needs as much care as the round function. If you ever design
a cipher, do not skimp on key mixing.
\end{trapbox}

\section{A full round, traced on bytes}
\label{ch12:trace}

Let us run one complete round on a concrete state. We will not do all ten rounds
by hand (that is the FIPS-197 appendix's job and runs for pages), but one round
shows how the four steps fit together. Start with the post-whitening state used in the
standard's example:

\[
\text{state} =
\begin{bmatrix}
\mathtt{19} & \mathtt{A0} & \mathtt{9A} & \mathtt{E9}\\
\mathtt{3D} & \mathtt{F4} & \mathtt{C6} & \mathtt{F8}\\
\mathtt{E3} & \mathtt{E2} & \mathtt{8D} & \mathtt{48}\\
\mathtt{BE} & \mathtt{2B} & \mathtt{2A} & \mathtt{08}
\end{bmatrix}.
\]

\textbf{Step 1, SubBytes.} Replace each byte by its S-box value. For example
$S(\mathtt{19})=\mathtt{D4}$, $S(\mathtt{A0})=\mathtt{E0}$,
$S(\mathtt{3D})=\mathtt{27}$, and so on. The state becomes

\[
\begin{bmatrix}
\mathtt{D4} & \mathtt{E0} & \mathtt{B8} & \mathtt{1E}\\
\mathtt{27} & \mathtt{BF} & \mathtt{B4} & \mathtt{41}\\
\mathtt{11} & \mathtt{98} & \mathtt{5D} & \mathtt{52}\\
\mathtt{AE} & \mathtt{F1} & \mathtt{E5} & \mathtt{30}
\end{bmatrix}.
\]

\textbf{Step 2, ShiftRows.} Rotate row $1$ left by $1$, row $2$ by $2$, row $3$
by $3$:

\[
\begin{bmatrix}
\mathtt{D4} & \mathtt{E0} & \mathtt{B8} & \mathtt{1E}\\
\mathtt{BF} & \mathtt{B4} & \mathtt{41} & \mathtt{27}\\
\mathtt{5D} & \mathtt{52} & \mathtt{11} & \mathtt{98}\\
\mathtt{30} & \mathtt{AE} & \mathtt{F1} & \mathtt{E5}
\end{bmatrix}.
\]

\textbf{Step 3, MixColumns.} Apply the matrix to each column. The first column is
$(\mathtt{D4},\mathtt{BF},\mathtt{5D},\mathtt{30})$ --- exactly the column we mixed
by hand earlier --- giving $(\mathtt{04},\mathtt{66},\mathtt{81},\mathtt{E5})$.
Mixing all four columns yields

\[
\begin{bmatrix}
\mathtt{04} & \mathtt{E0} & \mathtt{48} & \mathtt{28}\\
\mathtt{66} & \mathtt{CB} & \mathtt{F8} & \mathtt{06}\\
\mathtt{81} & \mathtt{19} & \mathtt{D3} & \mathtt{26}\\
\mathtt{E5} & \mathtt{9A} & \mathtt{7A} & \mathtt{4C}
\end{bmatrix}.
\]

\textbf{Step 4, AddRoundKey.} XOR in round key $1$, which for this example is

\[
\text{RoundKey}_1 =
\begin{bmatrix}
\mathtt{A0} & \mathtt{88} & \mathtt{23} & \mathtt{2A}\\
\mathtt{FA} & \mathtt{54} & \mathtt{A3} & \mathtt{6C}\\
\mathtt{FE} & \mathtt{2C} & \mathtt{39} & \mathtt{76}\\
\mathtt{17} & \mathtt{B1} & \mathtt{39} & \mathtt{05}
\end{bmatrix}.
\]

XORing cell by cell (e.g.\ $\mathtt{04}\xor\mathtt{A0}=\mathtt{A4}$) gives the
state that enters round $2$:

\[
\begin{bmatrix}
\mathtt{A4} & \mathtt{68} & \mathtt{6B} & \mathtt{02}\\
\mathtt{9C} & \mathtt{9F} & \mathtt{5B} & \mathtt{6A}\\
\mathtt{7F} & \mathtt{35} & \mathtt{EA} & \mathtt{50}\\
\mathtt{F2} & \mathtt{2B} & \mathtt{43} & \mathtt{49}
\end{bmatrix}.
\]

That is one round. Nine more like it (the last without MixColumns) and the block
is encrypted. Notice the avalanche already building: the single input byte
$\mathtt{19}$ in the top-left, after SubBytes, ShiftRows, and MixColumns, has influenced
an entire column of outputs, and the next round's ShiftRows will spread those
across all four columns. Figure~\ref{ch12:round-fig} summarizes the four layers
in order, and Figure~\ref{ch12:shiftrows-fig} shows the ShiftRows rotation that
carries bytes between columns.

\begin{figure}[H]
\centering
\begin{tikzpicture}[font=\small,
  box/.style={draw, rounded corners, minimum width=3.4cm, minimum height=0.85cm, align=center},
  sub/.style={box, fill=structurepurple!18},
  shiftr/.style={box, fill=sourceblue!18},
  mix/.style={box, fill=targetgreen!18},
  ark/.style={box, fill=quotientorange!22},
  arr/.style={-{Stealth[length=2.2mm]}, thick}]
  \node[box, fill=gray!12] (in) {state in};
  \node[sub, below=0.5cm of in] (sb) {\texttt{SubBytes}\\ \scriptsize confusion (nonlinear)};
  \node[shiftr, below=0.5cm of sb] (sr) {\texttt{ShiftRows}\\ \scriptsize move across columns};
  \node[mix, below=0.5cm of sr] (mc) {\texttt{MixColumns}\\ \scriptsize diffuse within columns};
  \node[ark, below=0.5cm of mc] (ak) {\texttt{AddRoundKey}\\ \scriptsize inject the secret};
  \node[box, fill=gray!12, below=0.5cm of ak] (out) {state out};
  \draw[arr] (in) -- (sb);
  \draw[arr] (sb) -- (sr);
  \draw[arr] (sr) -- (mc);
  \draw[arr] (mc) -- (ak);
  \draw[arr] (ak) -- (out);
  \node[right=0.9cm of mc, structurepurple, align=left] (key)
    {\scriptsize round key$_r$};
  \draw[arr, quotientorange] (key.west) .. controls +(left:0.5cm) and +(right:0.5cm) .. (ak.east);
\end{tikzpicture}
\caption{One AES round. The first three layers are fixed and public; only
\texttt{AddRoundKey} uses the secret. The final round omits \texttt{MixColumns}.}
\label{ch12:round-fig}
\end{figure}

\begin{figure}[H]
\centering
\begin{tikzpicture}[font=\scriptsize, scale=0.62, every node/.style={transform shape}]
  \foreach \r in {0,1,2,3} {
    \foreach \c in {0,1,2,3} {
      \node[draw, minimum size=0.8cm, fill=sourceblue!10] at (\c, -\r) {};
    }
  }
  \node at (1.5, 1.0) {\normalsize before};
  \node[draw, minimum size=0.8cm, fill=evered!25] at (0,-1) {$a_1$};
  \node[draw, minimum size=0.8cm, fill=evered!25] at (1,-1) {$a_5$};
  \node[draw, minimum size=0.8cm, fill=evered!25] at (2,-1) {$a_9$};
  \node[draw, minimum size=0.8cm, fill=evered!25] at (3,-1) {$a_{13}$};

  \begin{scope}[xshift=7cm]
  \foreach \r in {0,1,2,3} {
    \foreach \c in {0,1,2,3} {
      \node[draw, minimum size=0.8cm, fill=sourceblue!10] at (\c, -\r) {};
    }
  }
  \node at (1.5, 1.0) {\normalsize after (row $1$ rotated left by $1$)};
  \node[draw, minimum size=0.8cm, fill=evered!25] at (0,-1) {$a_5$};
  \node[draw, minimum size=0.8cm, fill=evered!25] at (1,-1) {$a_9$};
  \node[draw, minimum size=0.8cm, fill=evered!25] at (2,-1) {$a_{13}$};
  \node[draw, minimum size=0.8cm, fill=evered!25] at (3,-1) {$a_1$};
  \end{scope}
  \draw[-{Stealth[length=2mm]}, thick] (3.9,-1.5) -- (4.9,-1.5);
\end{tikzpicture}
\caption{\texttt{ShiftRows} on row $1$: a left cyclic rotation by one byte. Rows
$2$ and $3$ rotate by two and three. This is what carries bytes between columns so
the next \texttt{MixColumns} blends them.}
\label{ch12:shiftrows-fig}
\end{figure}

\section{Decryption: inverting each layer}

Because every layer is invertible, AES decryption simply applies the inverses in
reverse order. \texttt{InvSubBytes} uses the inverse S-box (apply the inverse
affine map, then field-invert). \texttt{InvShiftRows} rotates rows right instead
of left. \texttt{InvMixColumns} multiplies by the inverse of the MixColumns matrix
(whose entries are $\mathtt{0E},\mathtt{0B},\mathtt{0D},\mathtt{09}$ ---
deliberately larger, which is why naive decryption is a touch slower than
encryption). \texttt{AddRoundKey} is its own inverse.

With a small reordering trick (the ``equivalent inverse
cipher''), decryption can be written with the \emph{same} sequence of operations
as encryption, just with modified round keys, so hardware can share circuitry.
This symmetry is partly why the final round drops MixColumns --- it makes the
structure line up.

\begin{focusbox}
\textbf{AES is not a stream cipher and not its own inverse.} Encryption and
decryption use \emph{different} routines (inverse S-box, inverse matrix, reversed
key order). This differs from a Feistel cipher like DES, where decryption is
encryption with the round keys reversed and no separate inverse round function is
needed. The SPN structure trades that convenience for the ability to transform
every byte every round.
\end{focusbox}

\section{Why AES is secure}
\label{ch12:security}

``Secure'' for a block cipher means it behaves like a \emph{pseudorandom
permutation} (PRP): without the key, no efficient adversary can distinguish
$\Enc_k$ (for a random secret $k$) from a truly random permutation of
$128$-bit blocks, even given many chosen plaintext/ciphertext pairs. We cannot
prove AES is a PRP from any standard hardness assumption --- no one can prove that
of any practical cipher --- but we can show it resists every known general attack
with a comfortable margin, and that is the working standard of trust.

\begin{whybox}
\textbf{Confusion and diffusion in AES.} Claude Shannon's two design
goals are realized by specific layers.

\emph{Confusion} means the relationship between key and ciphertext is complicated
and nonlinear. SubBytes provides it: the only nonlinear step, with the best
available resistance to algebraic attacks.

\emph{Diffusion} means each plaintext (and key) bit influences many ciphertext
bits, so statistics smear out. ShiftRows + MixColumns provide it, as traced in
Figure~\ref{ch12:fig-wide-trail}. The \emph{wide-trail strategy} of Daemen and
Rijmen quantifies this: the MixColumns branch number is $5$, and over four rounds
any differential or linear trail must activate at least $25$ S-boxes. Each active
S-box contributes a factor bounding the trail's probability (at most $2^{-6}$ for
differentials), so any four-round trail has probability at most
$(2^{-6})^{25} = 2^{-150}$, far below the $2^{-128}$ threshold an attacker would
need. Ten rounds leave a large security margin.
\end{whybox}

\begin{figure}[htbp]
\centering
\begin{tikzpicture}[
  >=Stealth,
  cell/.style={draw=gray!55, fill=white, minimum size=3.4mm, inner sep=0pt, rounded corners=0.5pt},
  act/.style={draw=evered!85!black, fill=evered!35, minimum size=3.4mm, inner sep=0pt, rounded corners=0.5pt},
  arr/.style={->, thick, sourceblue!80!black},
  hdr/.style={font=\scriptsize\bfseries, text=gray!80!black, align=center}
]
  % Grid 1: 1 active byte
  \begin{scope}[shift={(0,0)}]
    \node[hdr] at (0.54, 0.72) {Start of Round $1$\\(\textnormal{1 active byte})};
    \foreach \r in {0,1,2,3} \foreach \c in {0,1,2,3} \node[cell] at (\c*0.36, -\r*0.36) {};
    \node[act] at (0, 0) {};
  \end{scope}

  \draw[arr] (1.35, -0.54) -- node[above, font=\scriptsize, align=center] {\texttt{SubBytes}\\+\,\texttt{ShiftRows}} (2.95, -0.54);

  % Grid 2: After SubBytes + ShiftRows (still 1 active byte in col 0)
  \begin{scope}[shift={(3.15,0)}]
    \node[hdr] at (0.54, 0.72) {Before \texttt{MixColumns}\\(\textnormal{1 active in col $0$})};
    \foreach \r in {0,1,2,3} \foreach \c in {0,1,2,3} \node[cell] at (\c*0.36, -\r*0.36) {};
    \node[act] at (0, 0) {};
  \end{scope}

  \draw[arr] (4.50, -0.54) -- node[above, font=\scriptsize, align=center] {\texttt{MixColumns}\\(branch $= 5$)} (6.10, -0.54);

  % Grid 3: End of Round 1 (4 active bytes in col 0)
  \begin{scope}[shift={(6.30,0)}]
    \node[hdr] at (0.54, 0.72) {End of Round $1$\\(\textnormal{4 active bytes})};
    \foreach \r in {0,1,2,3} \foreach \c in {0,1,2,3} \node[cell] at (\c*0.36, -\r*0.36) {};
    \foreach \r in {0,1,2,3} \node[act] at (0, -\r*0.36) {};
  \end{scope}

  \draw[arr] (7.65, -0.54) -- node[above, font=\scriptsize, align=center] {Round $2$\\ \texttt{ShiftRows}} (9.25, -0.54);

  % Grid 4: After Round 2 ShiftRows (1 active byte in each of the 4 columns)
  \begin{scope}[shift={(9.45,0)}]
    \node[hdr] at (0.54, 0.72) {After \texttt{ShiftRows}\\(\textnormal{1 per column})};
    \foreach \r in {0,1,2,3} \foreach \c in {0,1,2,3} \node[cell] at (\c*0.36, -\r*0.36) {};
    \node[act] at (0*0.36, -0*0.36) {};
    \node[act] at (3*0.36, -1*0.36) {};
    \node[act] at (2*0.36, -2*0.36) {};
    \node[act] at (1*0.36, -3*0.36) {};
  \end{scope}

  \draw[arr] (10.80, -0.54) -- node[above, font=\scriptsize, align=center] {Round $2$\\ \texttt{MixColumns}} (12.40, -0.54);

  % Grid 5: End of Round 2 (all 16 bytes active)
  \begin{scope}[shift={(12.60,0)}]
    \node[hdr] at (0.54, 0.72) {End of Round $2$\\(\textnormal{16 active bytes})};
    \foreach \r in {0,1,2,3} \foreach \c in {0,1,2,3} \node[act] at (\c*0.36, -\r*0.36) {};
  \end{scope}
\end{tikzpicture}
\caption{Two-round full diffusion via the wide-trail strategy. A single perturbed byte (red) spreads to all four bytes of its column after Round~1's \texttt{MixColumns} (branch number $5$). Round~2's \texttt{ShiftRows} distributes those four active bytes into four distinct columns, and Round~2's \texttt{MixColumns} expands each into a full column, activating all $16$ state bytes in just two rounds.}
\label{ch12:fig-wide-trail}
\end{figure}

\begin{attackbox}
\textbf{What attacks exist, and why they fall short.}

\emph{Brute force} on AES-128 needs $2^{128}$ trial decryptions --- physically
impossible (the energy alone exceeds what is available).

\emph{Differential cryptanalysis} tracks how input differences propagate. The
wide-trail bound above rules it out: no usable four-round differential exists,
let alone ten-round.

\emph{Linear cryptanalysis} looks for biased linear approximations; the same
$25$-active-S-box bound makes any biased relation vanish below detectability.

\emph{Biclique attack} (Bogdanov, Khovratovich, Rechberger, 2011) is the best
known key-recovery on full AES-128: it costs about $2^{126.1}$ --- a factor of
roughly $4$ better than brute force. After two decades of intense scrutiny, that
is the best known result against the full cipher.

\emph{Related-key attacks} on AES-192/256 (cost around $2^{99.5}$ for AES-256 in
the related-key model) are real but require the attacker to control key
differences, an assumption that almost never holds in practice.

The lesson: do not expect to break AES's math. The real attacks on AES deployments
are \emph{implementation} attacks --- below.
\end{attackbox}

\begin{attackbox}
\textbf{Cache-timing: the practical break.} The original software AES used a
$256$-entry table for SubBytes (and combined tables for the round). Table lookups
take different amounts of time depending on whether the indexed entry is in the
CPU cache, and the index depends on secret state. In 2005 Daniel Bernstein, and
independently Osvik--Shamir--Tromer, showed you can recover an AES key by
measuring timing or by priming and probing the cache from another process on the
same machine --- no break of the math required, just a leak through the
implementation. The fix that won: the \textbf{AES-NI} instruction set (Intel 2010,
now everywhere), which executes a whole round in hardware in constant time, with
no data-dependent table lookups. If you must do AES in software without AES-NI,
use a \emph{bitsliced} implementation that has no secret-indexed memory access.
\textbf{Never} ship the textbook table-lookup version in production.
\end{attackbox}

\begin{connectionbox}
The distinguishing game --- ``tell $\Enc_k$ apart from a random permutation'' ---
is the PRP security definition built in Chapters~8 and~11. The
$2^{-150}$ trail probabilities are union bounds over many possible trails; the
companion probability book's treatment of union bounds and tail estimates is the
exact tool. And the $2^{64}$ ``birthday'' limit on how many blocks you may safely
encrypt under one key in a chaining mode (next chapter) comes straight from the
birthday bound on a $128$-bit block size.
\end{connectionbox}

\section{Performance and where AES runs}

AES was designed to be fast in software \emph{and} cheap in hardware. The four layers are byte operations and XORs; the classic software
trick (Daemen and Rijmen's ``T-tables'') folds SubBytes, ShiftRows, and
MixColumns into four $1024$-byte lookup tables so a round is sixteen table reads
and a handful of XORs. With AES-NI, a single \texttt{aesenc} instruction does an
entire round; modern cores encrypt at several gigabytes per second per core, and
the constant-time hardware sidesteps the cache-timing problem entirely.

\begin{applicationbox}
\textbf{AES in the wild.}

\emph{TLS / HTTPS.} The dominant cipher suite is \texttt{AES-GCM}
(AES in counter mode with a Galois MAC), giving authenticated encryption. On a
phone without strong AES hardware, TLS 1.3 may instead pick ChaCha20-Poly1305,
which is faster in pure software --- but on anything with AES-NI, AES-GCM wins.

\emph{WiFi (WPA2 / WPA3).} The CCMP cipher is AES in CCM mode (counter mode plus
CBC-MAC). When you join a protected network, AES is encrypting your frames.

\emph{Disk encryption.} BitLocker, FileVault, LUKS, and Android/iOS file-based
encryption all use AES (in XTS mode for full-disk, where ciphertext stealing and
a tweak handle fixed-size sectors).

\emph{Signal / iMessage / Telegram.} The message payload encryption uses AES
(typically AES-CBC or AES-CTR with a separate HMAC, or AES-GCM).

\emph{VPNs and SSH.} IPsec, WireGuard's fallback paths, OpenVPN, and SSH all offer
AES-GCM or AES-CTR transport encryption.

In every case AES encrypts the bulk data; the asymmetric crypto (RSA, ECDH) only
sets up the AES key.
\end{applicationbox}

\begin{whybox}[Post-quantum margins and key sizes]
AES-128 offers $128$-bit security against classical attackers, far beyond the
reach of brute force. Grover's quantum
search algorithm would, in principle, halve the effective key strength to $64$
bits of quantum search --- which is why guidance for post-quantum margins often
recommends AES-256 (leaving $128$ bits of strength even against an idealized
quantum brute-forcer). Unlike RSA and elliptic curves, which Shor's algorithm
breaks outright, AES survives the quantum era essentially intact by just using a
longer key. This is one more reason AES is expected to stay in use for decades.
\end{whybox}

\section{Putting the whole cipher together}

For reference, here is the complete AES-128 encryption in pseudocode, every piece
of which we have now justified:

\begin{lstlisting}[language=Python]
# AES-128 encryption, structural pseudocode (LEARNING ONLY).
# Real use: call a vetted library (which uses AES-NI / constant time).
def aes128_encrypt(plaintext16, key16):
    state = to_state(plaintext16)          # 4x4 column-major bytes
    rk = key_expansion(key16)              # 11 round keys, words W0..W43

    add_round_key(state, rk[0])            # initial whitening
    for r in range(1, 10):                 # rounds 1..9
        sub_bytes(state)                   # nonlinear: S-box each byte
        shift_rows(state)                  # rotate rows 0,1,2,3 by 0,1,2,3
        mix_columns(state)                 # MDS mix within each column
        add_round_key(state, rk[r])        # XOR round key
    sub_bytes(state)                       # round 10: NO MixColumns
    shift_rows(state)
    add_round_key(state, rk[10])
    return from_state(state)
\end{lstlisting}

Each line maps to a design argument: SubBytes provides confusion, ShiftRows and
MixColumns provide diffusion, AddRoundKey injects the secret, the loop count sets
the security margin, and omitting the final MixColumns keeps decryption symmetric
at no security cost.

AES scrambles a single 16-byte block; to encrypt files, packets, and streams of arbitrary length without leaking patterns across blocks, Chapter~13 builds block-cipher modes of operation (ECB, CBC, CTR, CFB, and OFB).

\begin{recoverybox}
Before the exercises, a quick self-check. Can you (1) state the four round
operations and which one is nonlinear; (2) explain why ShiftRows and MixColumns
are useless apart but achieve full diffusion together; (3) derive an S-box entry from the field
inverse plus the affine map; (4) say why AES-128 uses ten rounds rather than,
say, four; and (5) name the real attack class that threatens AES deployments
(timing/cache side channels, not the math)? If any of these is shaky, reread the
matching section before moving on.
\end{recoverybox}

\section{Exercises}

\begin{exercise}
Compute $\texttt{xtime}(\mathtt{0x80})$, $\texttt{xtime}(\mathtt{0x7F})$, and
$\mathtt{0x02}\cdot\mathtt{0xFF}$ in $\F_{256}$ with the AES polynomial. Show the
shift-and-conditional-XOR for each.
\end{exercise}

\begin{exercise}
Verify by hand that $\mathtt{0x53}\cdot\mathtt{0xCA} = \mathtt{0x01}$ in
$\F_{256}$, confirming that $\mathtt{0xCA}$ is the multiplicative inverse of
$\mathtt{0x53}$ used in the worked S-box example. (Multiply the two polynomials
and reduce mod $m(x)$.)
\end{exercise}

\begin{exercise}
Using the affine formula $s_i = y_i \xor y_{i+4} \xor y_{i+5} \xor y_{i+6} \xor
y_{i+7} \xor c_i$ with $c=\mathtt{0x63}$, compute $S(\mathtt{0x00})$ and
$S(\mathtt{0x01})$ from their field inverses ($0^{-1}=0$ by convention,
$1^{-1}=1$). Confirm the famous values $S(\mathtt{0x00})=\mathtt{0x63}$ and
$S(\mathtt{0x01})=\mathtt{0x7C}$.
\end{exercise}

\begin{exercise}
\textbf{Diffusion.} Take any single column of four bytes, flip exactly one bit of
one byte, and run both versions through MixColumns. Show that all four output
bytes change. Explain how this demonstrates the branch number being greater than
$1$, and why ShiftRows is needed for the change to spread to other columns in the
next round.
\end{exercise}

\begin{exercise}
Run \texttt{ShiftRows} on the state whose rows are
$(\mathtt{00}, \allowbreak \mathtt{01}, \allowbreak \mathtt{02}, \allowbreak \mathtt{03})$,
$(\mathtt{10}, \allowbreak \mathtt{11}, \allowbreak \mathtt{12}, \allowbreak \mathtt{13})$,
$(\mathtt{20}, \allowbreak \mathtt{21}, \allowbreak \mathtt{22}, \allowbreak \mathtt{23})$,
$(\mathtt{30}, \allowbreak \mathtt{31}, \allowbreak \mathtt{32}, \allowbreak \mathtt{33})$. Write the resulting state.
\end{exercise}

\begin{exercise}
\textbf{Key schedule.} Continue the all-zero-key expansion from the worked
example: compute words $W_8$ through $W_{11}$ (so round key $2$). You will need
$\texttt{Rcon}_2 = \mathtt{02000000}$ and the S-box value
$S(\mathtt{0x63}) = \mathtt{0xFB}$. Show each step.
\end{exercise}

\begin{exercise}
Explain why AES omits \texttt{MixColumns} in its final round. Argue that if the
final MixColumns were present, an attacker could strip it off without knowing the
key, so it would add no security. (Hint: MixColumns is a fixed, public, invertible
linear map.)
\end{exercise}

\begin{exercise}
\textbf{Counting active S-boxes.} The wide-trail argument claims any four-round
differential trail activates at least $25$ S-boxes. Sketch the argument: one round
of MixColumns has branch number $5$, so a single active column in one round forces
at least $5$ active S-boxes across two consecutive rounds; combine across the
$4\times4$ structure over four rounds to reach $25$. Then estimate the trail
probability bound using $2^{-6}$ per active S-box.
\end{exercise}

\begin{exercise}
\textbf{Break a toy SPN.} Build a one-round ``AES'' with \emph{no} SubBytes (only
ShiftRows, MixColumns, AddRoundKey). Show that the entire map from plaintext to
ciphertext is affine over $\F_2$, and describe how an attacker recovers the round
key from a few known plaintext/ciphertext pairs by solving a linear system. This
is precisely why the nonlinear SubBytes layer cannot be removed.
\end{exercise}

\begin{exercise}
\textbf{Side channel.} Explain why a table-lookup implementation of SubBytes can
leak the key through cache timing, while an AES-NI hardware round or a bitsliced
software implementation does not. What property of the memory-access pattern is
the difference?
\end{exercise}

\begin{exercise}
AES-256 uses $14$ rounds and a slightly different key schedule than AES-128.
Given that the best single-key attack on AES-128 (biclique) is about $2^{126}$ and
related-key attacks on AES-256 reach about $2^{99.5}$, argue why ``more rounds and
a longer key'' did not straightforwardly make AES-256 stronger than AES-128 in the
related-key model, and why this is nonetheless not a practical concern for most
deployments.
\end{exercise}

\begin{exercise}
\textbf{Implementation.} Using the toy S-box generator from this chapter as a
start, write functions \texttt{shift\_rows}, \texttt{mix\_columns}, and
\texttt{add\_round\_key}, plus a \texttt{key\_expansion}, and reproduce the
single-round trace in Section~\ref{ch12:trace} (state in $\to$ state out). Then
encrypt the FIPS-197 test vector (plaintext
$\mathtt{00112233445566778899aabbccddeeff}$, key
$\mathtt{000102030405060708090a0b0c0d0e0f}$) and check you get ciphertext
$\mathtt{69c4e0d86a7b0430d8cdb78070b4c55a}$. State clearly in a comment that this
code is for learning, not production.
\end{exercise}
